A presente dissertação tem como objetivo o desenvolvimento de uma ferramenta computacional de apoio ao dimensionamento de elementos estruturais em betão armado, nomeadamente vigas, pilares e sapatas isoladas, em conformidade com o Eurocódigo 2.

Para esse efeito, foram formulados e implementados algoritmos de cálculo para cada um destes elementos estruturais, incorporando procedimentos iterativos que permitem obter soluções, construtivamente adequadas e otimizadas. Os modelos teóricos foram traduzidos para uma aplicação web desenvolvida em Python com recurso à framework Django, capaz de realizar o dimensionamento automático, gerar memórias de cálculo detalhadas e apresentar representações gráficas das soluções obtidas.

A fiabilidade da aplicação foi avaliada através de testes unitários automatizados e da comparação com exemplos de referência da literatura técnica, permitindo confirmar a coerência dos resultados produzidos. O trabalho desenvolvido demonstra, assim, a viabilidade da utilização de ferramentas computacionais no apoio ao dimensionamento estrutural, contribuindo para a digitalização e sistematização de procedimentos de cálculo em betão armado.

\vfill
\textbf{Palavras-chave:} Betão Armado; Eurocódigo 2; Dimensionamento Estrutural; Algoritmos; Python; Django; SVG; Validação Computacional.